\chapter{Conclusão}
\label{cap:conclusao}

% Chris Eng.ish --- V. Conclusion

\section{Síntese}

Este trabalho propôs e demonstrou, em escala de laboratório acadêmico, um método
e uma prova de conceito (ReqValLive) para validação live de Success Criteria:
aprovação do critério antes da run, avaliação de telemetria MQTT real, laudo
PASS/FAIL e retorno estruturado ao modelo SysML~v2 (SysON).

\section{Contribuições}

\begin{enumerate}
  \item \textbf{Metodológica:} operacionalização de gate e evidência temporal
        para SC live; vocabulário de naturezas de observação (GSE).
  \item \textbf{Tecnológica:} artefacto ReqValLive + fecho SysON.
  \item \textbf{Empírica:} demonstração no CONCEPTIO com drones reais / MQTT.
\end{enumerate}

\section{Limitações}

Escala lab; cobertura parcial das naturezas de observação; SysML como esqueleto /
write-back, não simulação Magic; SARITA não é pré-requisito de conclusão.

\section{Trabalhos futuros}

\begin{itemize}
  \item Estender o motor live (taxa, tempo/dwell, eventos, estados) conforme
        taxonomia GSE e telemetria disponível.
  \item UI e missão MQTT alinhadas ao padrão do orientador (RPS).
  \item Evolução TWC REST / plugin Magic em paralelo, sem bloquear a escrita.
  \item Ensaio SARITA/CET-ADS se o evento ocorrer.
\end{itemize}
